% Intended LaTeX compiler: pdflatex
\documentclass[a4paper,12pt]{article}
\usepackage[utf8]{inputenc}
\usepackage[T1]{fontenc}
\usepackage{graphicx}
\usepackage{longtable}
\usepackage{wrapfig}
\usepackage{rotating}
\usepackage[normalem]{ulem}
\usepackage{amsmath}
\usepackage{amssymb}
\usepackage{capt-of}
\usepackage{hyperref}
\usepackage{\string~/.emacs.d/common}
\author{mahmood}
\date{\today}
\title{common lisp}
\hypersetup{
 pdfauthor={mahmood},
 pdftitle={common lisp},
 pdfkeywords={},
 pdfsubject={},
 pdfcreator={Emacs 28.2 (Org mode 9.5.5)}, 
 pdflang={English}}
\begin{document}

\maketitle
\begin{note}
this document and others can be found on my website\\
\url{https://mahmoodsheikh36.github.io/}
\end{note}
\section{for loop}
\label{sec:org485014a}
\lstset{language=Lisp,label= ,caption= ,captionpos=b,numbers=none}
\begin{lstlisting}
(loop for x in '(1 2 3)
      collect x) ;; returns list
(loop for x in '(1 2 3)
      do (print x)) ;; prints the numbers
\end{lstlisting}
\section{functions}
\label{sec:org1e257b2}
\lstset{language=Lisp,label= ,caption= ,captionpos=b,numbers=none}
\begin{lstlisting}
(defun <name> (list of arguments)
  "docstring"
  (body))
\end{lstlisting}
\section{plist}
\label{sec:org918814f}
defining a plist\\
\lstset{language=Lisp,label= ,caption= ,captionpos=b,numbers=none}
\begin{lstlisting}
(setf my-plist '(alt "yes" bar "no"))
\end{lstlisting}
retrieve a value by key\\
\lstset{language=Lisp,label= ,caption= ,captionpos=b,numbers=none}
\begin{lstlisting}
(getf my-plist 'alt)
\end{lstlisting}
change the value that corresponds to a key\\
\lstset{language=Lisp,label= ,caption= ,captionpos=b,numbers=none}
\begin{lstlisting}
(setf (getf my-plist 'alt) 45)
\end{lstlisting}
get a list of the attributes\\
\lstset{language=Lisp,label= ,caption= ,captionpos=b,numbers=none}
\begin{lstlisting}
(loop for (key value) on my-plist by #'cddr
      collect key)
\end{lstlisting}
\section{unbind variable}
\label{sec:orge636551}
\lstset{language=Lisp,label= ,caption= ,captionpos=b,numbers=none}
\begin{lstlisting}
(makunbound 'foo)
\end{lstlisting}
\section{undefine funciton}
\label{sec:org217abd2}
\lstset{language=Lisp,label= ,caption= ,captionpos=b,numbers=none}
\begin{lstlisting}
(fmakunbound 'function)
\end{lstlisting}
\section{quicklisp}
\label{sec:org4d7902c}
need to install it manually afaik\\
\lstset{language=bash,label= ,caption= ,captionpos=b,numbers=none}
\begin{lstlisting}
curl -O https://beta.quicklisp.org/quicklisp.lisp
sbcl --load quicklisp.lisp
\end{lstlisting}
then in sbcl\\
\lstset{language=Lisp,label= ,caption= ,captionpos=b,numbers=none}
\begin{lstlisting}
(quicklisp-quickstart:install :path "~/.quicklisp")
(ql:add-to-init-file)
\end{lstlisting}
then to install a package in sbcl:\\
\lstset{language=Lisp,label= ,caption= ,captionpos=b,numbers=none}
\begin{lstlisting}
(ql:quickload "cl-csv")
\end{lstlisting}
\section{classes}
\label{sec:orgef1aa29}
example class:\\
\lstset{language=Lisp,label= ,caption= ,captionpos=b,numbers=none}
\begin{lstlisting}
(defclass my-class ()
  ((property1 :accessor property1-accessor)
   (property2 :accessor property2-accessor)))
\end{lstlisting}

the \texttt{accessor} is the name you use to access a property1\\
we construct an object using \texttt{make-instance}:\\
\lstset{language=Lisp,label= ,caption= ,captionpos=b,numbers=none}
\begin{lstlisting}
(defvar p1 (make-instance 'my-class))
\end{lstlisting}

we can initialize the class with specific values, first we need to use the \texttt{:initarg} keyword on the class properties\\
\lstset{language=Lisp,label= ,caption= ,captionpos=b,numbers=none}
\begin{lstlisting}
(defclass my-class ()
  ((property1 :initarg property1-initarg :initform default-value :accessor property1-accessor)
   (property2 :accessor property2-accessor)))
\end{lstlisting}

then we can pass the values we want to \texttt{make-instance}\\
\texttt{:initform} is used to pass a default value to the slot incase \texttt{:initarg} argument isnt passed to \texttt{make-instance}\\
\lstset{language=Lisp,label= ,caption= ,captionpos=b,numbers=none}
\begin{lstlisting}
(defvar p1 (make-instance 'my-class :property1-initarg property1-value))
\end{lstlisting}

this is the syntax for accessing variables of a class\\
\lstset{language=Lisp,label= ,caption= ,captionpos=b,numbers=none}
\begin{lstlisting}
(property1-accessor p1)
\end{lstlisting}
we can do inheritance by passing the parent class in the parenthases after the class name:\\
\lstset{language=Lisp,label= ,caption= ,captionpos=b,numbers=none}
\begin{lstlisting}
(defclass my-class (my-parent-class))
\end{lstlisting}
the following syntax is used to define a function for a class, note that the function \texttt{print-object} is the function that gets called when trying to print an object using \texttt{(print object)} so thats what we're defining here\\
\lstset{language=Lisp,label= ,caption= ,captionpos=b,numbers=none}
\begin{lstlisting}
(defmethod print-object ((obj my-class) stream)
  (print-unreadable-object (obj stream :type t)
    (format stream "value of property1: ~a" (property1 obj))))
\end{lstlisting}
\section{structs}
\label{sec:org70917ce}
structs are basically simpler classes with default functions/initializers/etc\\
\lstset{language=Lisp,label= ,caption= ,captionpos=b,numbers=none}
\begin{lstlisting}
(defstruct rgb ()
           (r 0) (g 0) (b 0)) ;; initialize slot values to 0

;; e.g. constructor built by default
(setf color (make-rgb :r 2)) ;; => #S(RGB :NIL NIL :R 2 :G 0 :B 0)

(list (rgb-r color) (rgb-b color)) ;; => (2 0)
\end{lstlisting}
\section{arrays}
\label{sec:org5f6c6ac}
to create an array:\\
\lstset{language=Lisp,label= ,caption= ,captionpos=b,numbers=none}
\begin{lstlisting}
(make-array 10 :initial-element 10)
\end{lstlisting}

\begin{tabular}{rrrrrrrrrr}
10 & 10 & 10 & 10 & 10 & 10 & 10 & 10 & 10 & 10\\
\end{tabular}

multidimensional arrays:\\
\lstset{language=Lisp,label= ,caption= ,captionpos=b,numbers=none}
\begin{lstlisting}
(make-array '(4 3) :initial-element 15)
\end{lstlisting}

\begin{verbatim}
#2A((15 15 15) (15 15 15) (15 15 15) (15 15 15))
\end{verbatim}


access slot in array\\
\lstset{language=Lisp,label= ,caption= ,captionpos=b,numbers=none}
\begin{lstlisting}
(aref (make-array '(4 3 1) :initial-element 15) 2 0 0)
\end{lstlisting}

\begin{verbatim}
15
\end{verbatim}


get the length of the array (last argument is the dimension to check length of):\\
\lstset{language=Lisp,label= ,caption= ,captionpos=b,numbers=none}
\begin{lstlisting}
(print (array-dimension (make-array '(4 3) :initial-element 15) 0))
(print (array-dimension (make-array '(4 3) :initial-element 15) 1))
\end{lstlisting}

\begin{verbatim}

4 
3 
\end{verbatim}


iterate through array:\\
\lstset{language=Lisp,label= ,caption= ,captionpos=b,numbers=none}
\begin{lstlisting}
(map nil #'print (make-array 3 :initial-element 15))
\end{lstlisting}

\begin{verbatim}

15 
15 
15 
\end{verbatim}


iterate with index:\\
\lstset{language=Lisp,label= ,caption= ,captionpos=b,numbers=none}
\begin{lstlisting}
(let ((arr (make-array '(4 3) :initial-element 15)))
  (loop for i from 0 below (array-dimension arr 0)
    do (print (aref arr i 1))))
\end{lstlisting}

\begin{verbatim}

15 
15 
15 
15 
\end{verbatim}


turn multidimensional list into array:\\
\lstset{language=Lisp,label= ,caption= ,captionpos=b,numbers=none}
\begin{lstlisting}
(defun list-dimensions (list depth)
  (loop repeat depth
        collect (length list)
        do (setf list (car list))))
(defun list-to-array (list depth)
  (make-array (list-dimensions list depth)
              :initial-contents list))
\end{lstlisting}

use displacement arrays to get the first element of every row in a 2d array:\\
\lstset{language=Lisp,label= ,caption= ,captionpos=b,numbers=none}
\begin{lstlisting}
(loop for i from 0 below (array-dimension arr 0)
      collect (let ((row (make-array
                          (array-dimension arr 1)
                          :displaced-to arr
                          :displaced-index-offset (* i (array-dimension arr 1)))))
                (aref row 0)))
\end{lstlisting}

we can create resizable arrays with the \texttt{:adjustable} keyword, \texttt{:fill-pointer} denotes the next position to be filled in the array and is automatically maintained by the array itself, but has to be initialized\\
\lstset{language=Lisp,label= ,caption= ,captionpos=b,numbers=none}
\begin{lstlisting}
(defparameter *x* (make-array 0 :adjustable t :fill-pointer 0))
(vector-push-extend 'a *x*)
(vector-push-extend 'b *x*)
*x*
\end{lstlisting}

\begin{tabular}{ll}
A & B\\
\end{tabular}

\section{vector}
\label{sec:org1734a31}
create vector\\
\lstset{language=Lisp,label= ,caption= ,captionpos=b,numbers=none}
\begin{lstlisting}
(vector 3 5 1)
\end{lstlisting}

we can use \texttt{length} on vectors\\
\lstset{language=Lisp,label= ,caption= ,captionpos=b,numbers=none}
\begin{lstlisting}
(length (vector 3 5 1)) ==> 3
\end{lstlisting}

get nth element:\\
\lstset{language=Lisp,label= ,caption= ,captionpos=b,numbers=none}
\begin{lstlisting}
(defparameter *x* (vector 1 2 3))
(length *x*) ==> 3
(elt *x* 1)  ==> 2
\end{lstlisting}

we can have nested vectors\\
\lstset{language=Lisp,label= ,caption= ,captionpos=b,numbers=none}
\begin{lstlisting}
(let ((a (vector (vector 1 2) (vector 3) 4)))
  (print a)
  (print (elt a 0)))
\end{lstlisting}

\begin{verbatim}

#(#(1 2) #(3) 4) 
#(1 2) 
\end{verbatim}


converting nested lists to nested vectors:\\
\lstset{language=Lisp,label= ,caption= ,captionpos=b,numbers=none}
\begin{lstlisting}
(defun list->vector (mylist)
  "convert nested lists into nested vectors"
  (coerce
   (loop for sublist in mylist
         collect
         (cond
           ((equal (type-of sublist) 'cons) (list->vector sublist))
           (t sublist)))
   'vector))
(let ((mylist '((1 2) 3 4)))
  (print (list->vector mylist))) ;; => #(#(1 2) 3 4) 
\end{lstlisting}

someone on discord gave me this shorter one though:\\
\lstset{language=Lisp,label= ,caption= ,captionpos=b,numbers=none}
\begin{lstlisting}
(defun list->vector (list)
  (if (atom list) list
      (map 'vector #'list->vector list)))
\end{lstlisting}
\section{multithreading}
\label{sec:org966b94d}
\texttt{sb-thread:make-thread} takes a function to call in a newly created thread.\\
\lstset{language=Lisp,label= ,caption= ,captionpos=b,numbers=none}
\begin{lstlisting}
(sb-thread:make-thread
 (lambda ()
   (progn
     (sleep 0) ;; give other threads a chance to run, then return here
     (setf c (+ a b))
     (print "ADDITION:")
     (print c))))
\end{lstlisting}
\section{resources}
\label{sec:org7d64a87}
\url{https://lisp-journey.gitlab.io/blog/common-lisp-macros-by-example-tutorial}\\
\url{https://taeric.github.io/CodeAsData.html}\\
\url{https://github.com/ghollisjr/cl-ana}\\
\url{https://github.com/marcoheisig/Petalisp}
\end{document}